\documentclass[11pt,a4paper]{article}
\usepackage[utf8]{inputenc}
\usepackage[T1]{fontenc}
\usepackage{amsfonts}
\usepackage[french]{babel}
\frenchbsetup{StandardLists=true} % à inclure si on utilise \usepackage[francais]{babel}
\usepackage{enumitem}
\usepackage{amssymb}
\usepackage{amsmath}
\usepackage[left=2cm,right=2cm,top=2cm,bottom=2cm]{geometry}
\usepackage{multirow}
\usepackage[dvipsnames]{xcolor}

\usepackage[T1]{fontenc}
\usepackage{fouriernc}
\usepackage{graphicx}

\usepackage{setspace}
\singlespacing

\usepackage{pstricks}
\usepackage{xkeyval}

\usepackage{pst-labo}
\usepackage{chemist} %chimie


\usepackage{multicol}
\usepackage{caption}
\usepackage{fancyhdr}
\pagestyle{fancy}
\renewcommand\headrulewidth{1pt}
\fancyhead[L]{Lycée Denis Diderot }
\fancyhead[R]{Classe de 1 Bac Pro }
\setcounter{page}{1}\fancyfoot[C]{page \thepage}
\fancyfoot[R]{Année 2024-2025}
\fancyfoot[L]{Alexis RUIZ}
\usepackage{multirow,tabularx,booktabs}
\usepackage[tikz]{bclogo} 

\newcommand{\cadreblanc}[1]{%
\medskip\framebox[\linewidth][c]{%
\rule{0mm}{#1mm}}\par
\medskip}

\usepackage{awesomebox}
\setlength{\aweboxvskip}{0mm}
\usepackage{pifont}
\usepackage{chemmacros}
\chemsetup{modules=all}
\setlength{\columnseprule}{.25mm}

\renewcommand{\arraystretch}{2.5}
\newcommand*\acocher{\quitvmode{\fboxsep0pt \fboxrule0.8pt \fbox{\vrule height1.8ex depth.1ex width0pt \vrule height0pt depth0pt width1.9ex }}}

\newcolumntype{C}[1]{>{\centering\arraybackslash }b{#1}}

\newcommand{\code}{\awesomebox[black]{2pt}{

\includegraphics[scale=0.15]{Chapitre 1 - Energie et puissance électrique/images/frame.png}
}{black}}

\newcommand{\moteur}{\awesomebox[black]{2pt}{

\includegraphics[scale=1]{images/moteur.jpg} 
}{black}}

\newcommand{\defconclusion}{\awesomebox[black]{2pt}{\rotatebox{90}{\Large{\textbf{Conclusion}}}}{black}}

\newcommand{\defbox}{\awesomebox[black]{2pt}{
\includegraphics[scale=0.5]{Chapitre 1 - Energie et puissance électrique/images/appel exam.png}}{black}}

  \newenvironment{itemize*}%  % listes compactes
    {\begin{itemize}%
      \setlength{\itemsep}{0pt}%
      \setlength{\parskip}{0pt}}%
    {\end{itemize}}

\frenchbsetup{StandardLists=true}



\begin{document}

\flushleft \textbf{NOM : \hfill Classe : \hspace{3cm} ~~\\[2mm]
Prénom:\\[0,4cm]
}
\hrule\
\\[0,2cm]

\begin{center}
 \LARGE{TP2 - Préparation d'une boisson isotonique.    }
\end{center}

\vspace{3 mm}

\normalsize

\hrule\
\\[0,2cm]
   \begin{center}
   \textbf{Le soin et la rédaction interviendront dans la notation.}
   \end{center} 

   Vous allez préparer 100 mL d'une solution dont la concentration en sucre de saccharose \ch{C_12H_22 O_11} est \ch{C_1=} 0,10 mol/L et celle en sel de cuisine (chlorure de sodium, NaCl) est \ch{C_2=}0,020 mol/L.
\noindent%
\begin{minipage}{.35\textwidth}%

    
 \begin{bclogo}[couleur=white,  arrondi =0.1 , logo = \bcinfo, epBarre =2]{ Matériel }

\begin{itemize*}
    \item 1 balance
    \item 2 coupelles
    \item 1 fiole jaugée de 100 mL
    \item 1 entonnoir
\end{itemize*}
\end{bclogo}

    
 \begin{bclogo}[couleur=white,  arrondi =0.1 , logo = \bcinfo, epBarre =2]{ Produits }

\begin{itemize*}
    \item 50 g de sucre blanc
    \item Sel de table
    \item Pissette eau du robinet
\end{itemize*}
\end{bclogo}


\end{minipage}%
\hfill
\begin{minipage}{.60\textwidth}%
\fbox{\includegraphics[width=\textwidth]{Chapitre 4 - Solutions aqueuses/images/TP2.png}}
\end{minipage}%



\fcolorbox{black}{blue!30}{\textbf{Calculs préparatoires}}
\begin{itemize*}
    \item Calculer, en g/mol, la masse molaire M du saccharose de formule  \ch{C_12H_220_11}.
    
    \cadreblanc{15}
\vfill 

    \item En utilisant la formule du cours liant la concentration \textbf{c} en mol/L , le nombre de mole \textbf{n} et le volume \textbf{V}, $c=\dfrac{n}{V}$ , calculer , en mol, \textbf{n}, la quantité de matière de saccharose nécessaire à la préparation de la solution. 

    \cadreblanc{15}
    
    \vfill 
    
    \item Calculez la masse m, en g, de saccharose nécessaire à la préparation de cette solution, en utilisant la formule $n=\dfrac{m}{M}$. 
    
    \cadreblanc{15}

    \vfill 
    
        \item Calculez la masse m, en g, de chlorure de sodium nécessaire à la préparation de cette solution. 
    \cadreblanc{35}
    
    \defbox{\fbox{
\begin{minipage}{8cm}{

Appeler M Ruiz pour qu'il vérifie vos calcul.}
\end{minipage}
}}

\vspace{3mm}
\fcolorbox{black}{blue!30}{\textbf{Protocole}}

\begin{itemize}[label=$\bullet$]
    \item Dans une coupelle, pesez 3,42g de sucre de saccharose en poudre.

    \vspace{3mm}

    \item Dans l’autre coupelle,pesez 0,12g de chlorure de sodium. 

    \vspace{3mm}

    \item Introduisez le sucre de saccharose puis le chlorure de sodium dans la fiole jaugée à l’aide de l’entonnoir.

    \vspace{3mm}

    \item Remplissez-la aux trois-quarts avec l’eau de la pissette.

    \vspace{3mm}
    
    \item Bouchez et agitez honrizontalement pour dissoudre. 

    \vspace{3mm}
    
    \item Ajoutez de l’eau jusqu’au trait de jauge. Rebouchez et agitez.
    
    
\end{itemize}



  \defbox{\fbox{
\begin{minipage}{8cm}{

Appeler M Ruiz pour qu'il vérifie votre solution.}
\end{minipage}
}}



\end{itemize*}

\dotfill 

\vfill 

\begin{center}
    


\fbox{\includegraphics[scale=0.74]{Chapitre 4 - Solutions aqueuses/images/Dilution.png}}
\end{center}

\vfill 

\end{document}